\section{The actual reader on its small panel}
\label{sec:ex-lattice-actual-reader-small-window}

Fix the outer source scale $x$ and let $B$ denote its outer radius.
The sector panel radius is the natural quotient
\[
 p=\left\lfloor\frac{B+80}{20}\right\rfloor.
\]
The matrices, forest, frozen coordinates, masks, and cutoff below are
always those of the outer-$B$ construction.  The smaller $p$ is the
panel already chosen by that construction; it is not a new source scale.
Write $\mathcal D_U(R)$ for the seed-dome of radius $R$ in the fixed
configuration $U$.  Outer clearance means
$o\notin\mathcal D_U(11B+1)$.

\begin{lemma}[Two-scale containment]
\label{lem:ex-lattice-small-window-geometry}
If $B\ge200$, then
\begin{equation}\label{eq:ex-lattice-small-window-geometry}
\begin{aligned}
 p&\le B/20+4, & 30p&<11B+1-p,\\
 9p&\le B, & B/4+p+B/50&\le B/2,\\
 B/4+2p+B/25&\le B, &p+B/50+B/4+B/2&\le B.
\end{aligned}
\end{equation}
Thus a radius-$8p$ ball centered at any panel center within distance
$p$ of $o$ lies in the outer radius-$B$ ball.  A target within $B/4$
and its row span of at most $B/50$ lie inside the radius-$B/2$ region
where the source-row bound is available.  Adding the further dependency
radius $B/2$ still stays inside the outer $B$ neighborhood.  Every
panel center moved by at most $p$ remains farther than $30p$ from each
seed.
\end{lemma}

\begin{proof}
The source uses natural division, so $p\le(B+80)/20=B/20+4$.
The seed-dome radius is $10B+B+1=11B+1$.  Moving a center by at most
$p$ leaves clearance at least $11B+1-p$, and
\[
 11B+1-p-30p\ge 11B+1-31(B+80)/20
                  =(189B-2460)/20>0.
\]
Also $9p\le9(B+80)/20\le B$ for $B\ge200$.  The total-reach bound
follows from
\[
 p+B/50+B/4+B/2\le B/20+4+ B/50+B/4+B/2
                         =(41B+200)/50\le B.
\]
Finally, $B/4+p+B/50\le B/2$ for $B\ge200$.  These inequalities
and $B/4+2p+B/25\le B$ follow on substituting $p\le B/20+4$.
They concern only support and source queries; no matrix, cutoff or mask
is changed.
\end{proof}

\begin{proposition}[Actual local mean and panel covariance]
\label{prop:ex-lattice-actual-reader-panel}
Fix $R\ge1$ and put
\[
 \chi(x)=16\{2(B+1)+1\}^5n_{\rm plaq}\,\vartheta(x).
\]
Here $n_{\rm plaq}$ is the finite number of plaquette orientations and
$\vartheta(x)$ is the source temperature, and
$B=\lceil(\log x)^2\rceil$. There are a fixed finite
number of row, color, and local jet types, and a radius-$r$ ball in the
four-dimensional lattice has at most $(2r+1)^4$ sites. Consequently
the row and panel counts below are bounded by a fixed multiple of
$(1+B)^4$.
Fix one patch, tube, forest, tag, posed configuration and history over
the entire panel.  Assume the source's branch/funding, floor,
attachment, swornness, diagonalization and roll conditions, and the
same-pose hypotheses at outer radius $B$: for a source threshold
$\delta>0$, the coarse side satisfies $c\ge100(1+B)$, and the stored
flags equal the flags of the posed configuration at temperature
$\vartheta(x)$. For each $x$ under consideration, fix $t$ with
$0\le t\le\min\{\delta,R\chi(x)\}$.
Assume that the direction $v$ has size at most $t$, the local price bound
holds on that ball, its group representative agrees there with the posed
field, and the center is outside $\mathcal D_U(11B+1)$. Every shifted center used in the proof also satisfies the
same pose margin: its distance from the fixed pose center plus the
pose's wall allowance is at most the original width $W$.  Use the
actual panels selected by the outer construction and the
official-address map to the cleared panel. The mean conclusion is for
an active coordinate $i$ whose anchor is within $B/4$ of $o$; the
column conclusion is for every $j$ in that same selected panel.

Let $\eta>0$ be the minimum decay rate in the short-kernel approximation
and long-path budget. Put
\begin{equation}\label{eq:ex-lattice-reader-widening}
 \operatorname{widen}_x(t)=(1+\operatorname{till}(x))t
                         +e^{-\eta\operatorname{till}(x)}.
\end{equation}
Since $t\le R\chi(x)$ and $\chi(x)\to0$, the source small-price
threshold holds for all sufficiently large $x$.
Let $D_x$ be the number of source rows that can contribute to a fixed
normal-residual coordinate; the four-dimensional row count gives
$D_x\le C(1+B)^4$. Let $n_p$ bound the number of coordinates in a
reader panel at radius $p$. There are finitely many color and local jet
types, so the four-dimensional panel count gives
$n_p\le C(1+p)^4\le C(1+B)^4$. The source assumptions
include its proved fixed-row chart and path bounds, its actual and bare
positive Gram gaps $\lambda,\Lambda>0$, their uniform row constants,
and the finite inverse-window estimates from the canonical local
covariance argument.  Then, for all sufficiently large $x$, uniformly over these admissible $t$,
the period and this same history,
\begin{align}
 |\mu_i-a_i|&\le x^{-7/8},
 \label{eq:ex-lattice-reader-mean}\\
 |((Q_V|_{W_i})^{-1})_{ji}-((Q_1|_{W_i^0})^{-1})_{\bar j\bar i}|
      &\le x^{-1/4}.
 \label{eq:ex-lattice-reader-column}
\end{align}
Here $a_i$ is the frozen stationary-response coordinate, $\mu_i$ is
the actual local Gaussian mean, $Q_V$ and $Q_1$ are respectively the
actual and identity-configuration row Gram matrices, $W_i$ is the
actual reader panel, $W_i^0$ is its cleared counterpart, and bars
denote the official-address bijection.  The matrices are not asserted
to have equal entries.  The error in the second line is the local
actual-to-cleared inverse-column comparison; the later covariance
split that upgrades it to the full physical-tensor estimate is not part
of this proposition.
\end{proposition}

\begin{proof}
\emph{The nonlinear row error.}
Set $m$ to be the frozen stationary response to $v$.  The exact
first-order identity for each uncut row is
\begin{equation}\label{eq:ex-lattice-frozen-row-cancel}
 I_r(m)+F_r(v)=0,
\end{equation}
where $F_r(v)$ is the linear first variation and $I_r(m)$ is the
corresponding flat pulse. Write $B_r(V)$ for the smooth base
contribution and $P_r(V,m)$ for the marked nonlinear pulse. The
selected row, after subtraction of its stationary linear part, is
exactly
\[
 \Delta_r=s\bigl[(B_r(V)-F_r(v))+(P_r(V,m)-I_r(m))\bigr].
\]
To verify this row formula, restrict to the finite read stencil.
Outer clearance makes every free bond there active, while the frozen
response $m$ vanishes on tree bonds. Thus the marked direction agrees
with $m$ throughout the stencil. The additional frozen-input field
also vanishes there, so finite row dependence removes it from the
selected row formula. No value outside the stencil is used.
Indeed, the selected row is $s(B_r(V)+P_r(V,m))$, while
\eqref{eq:ex-lattice-frozen-row-cancel} says $F_r(v)+I_r(m)=0$.
Unselected rows vanish by the row mask.  This cancellation is made
before taking absolute values.

For the base difference, write $F_j$ for the actual transport factor
at the $j$th position of a path and $H_j$ for its frozen linear
factor. Expand their finite transport products by the exact
telescoping identity
\[
 F_1\cdots F_m-H_1\cdots H_m
 =\sum_{j=1}^mF_1\cdots F_{j-1}(F_j-H_j)H_{j+1}\cdots H_m.
\]
The local chart bounds each changed factor by the field size times its
path length; summing the retained lengths gives
$(1+\operatorname{till}(x))t$. The omitted paths contribute at most
$e^{-\tau\operatorname{till}(x)}$. Thus
$|B_r(V)-F_r(v)|\le A t\,\operatorname{widen}_x(t)$.
For the pulse difference, first restrict the direction to the outer
comparison ball of radius $B$. Finite row dependence makes the row
insensitive to the discarded part. The pulse construction then cuts
paths at $\operatorname{till}(x)$: its finite pulse bound is
$Z(1+\operatorname{till}(x))t\|m\|$ and the weighted omitted-path
bound is $Z e^{-\sigma\operatorname{till}(x)}\|m\|$. The carry
estimate gives $\|m\|\le Xt$. Consequently, with
$\eta=\min\{\tau,\sigma\}$,
$|P_r(V,m)-I_r(m)|\le ZXt\,\operatorname{widen}_x(t)$.
Combining these inequalities proves, for a constant $A_0$,
\begin{equation}\label{eq:ex-lattice-row-delta}
 |\Delta_r|\le A_0t\,\operatorname{widen}_x(t),
 \qquad d(o,r)\le B/2.
\end{equation}
This derivation uses the same frozen response and the same outer
cutoff for every row; it does not reset the gauge or history panel by
panel.

\emph{Normal residual and the mean.}
Let $M$ denote the actual row operator, $\Pi_0$ the projection onto
inactive coordinates, $A=M^*M+\Pi_0$, and let $r_0$ be the affine
source vector with $\ell=M^*r_0$.  Let $\mathsf P$ be the padded
principal inverse used by the local Gaussian mean. For the frozen
candidate $a$, the selected affine residual is $r(a)=Ma+r_0$; on each
selected row its coordinate is the change $\Delta_r$ estimated above.
The exact normal
residual identities are
\begin{equation}\label{eq:ex-lattice-normal-residual}
 \mu=-\mathsf P^*\ell,\qquad
 \mu-a=E^*a-\mathsf P^*(Aa+\ell),\qquad
 Aa+\ell=M^*r(a)+\Pi_0a,
 \quad E=A\mathsf P-I.
\end{equation}
These are the same finite-dimensional identities as
\eqref{eq:mean-error-identity}. The candidate $a$ is extended by zero
outside the active coordinates, so $\Pi_0a=0$. All queried mean
coordinates are active; no assertion that every globally inactive
coordinate is a tree coordinate is used.
Each contributing row entry of $M$ is bounded by a fixed constant
$M_0$, and at most $D_x$ rows contribute.  Equation
\eqref{eq:ex-lattice-row-delta} therefore gives
\[
 |(Aa+\ell)_j|\le D_xM_0A_0t\,\operatorname{widen}_x(t).
\]
The local inverse-column sum is bounded by $N_0$. The inverse-window
estimate is used only at the queried central coordinate:
$|(E^*a)_i|\le E_0e^{-sB}Xt$. No global operator-norm estimate for
$E$ is asserted. Since $\|a\|_\infty\le Xt$,
\eqref{eq:ex-lattice-normal-residual} implies
\begin{equation}\label{eq:ex-lattice-focused-mean-bound}
 |\mu_i-a_i|
 \le N_0D_xM_0A_0t\,\operatorname{widen}_x(t)+E_0e^{-sB}Xt.
\end{equation}
In the row-Gram formulation used for the actual reader, let $Q$ be
the row Gram matrix, $b$ its affine source, and $\nu_W=Q[W]^{-1}$
extended by zero off the selected panel.  The corresponding exact
identity is
\[
 -\nu_W^*b-a=(Q\nu_W-I)^*a-\nu_W^*(Qa+b),
\]
This form makes the local residual sum and the off-window tail
separate.

For the source payment $t\le R\chi(x)$. Thus
$t\le C_R(1+B)^5\vartheta(x)$. Put
$b_x=1+B$ and $\mathcal B_R(x)=16R3^5n_{\rm plaq}b_x^5$; then
$R\chi(x)\le\mathcal B_R(x)\vartheta(x)$. With
$u=x^{-1/2+1/32}$, the temperature and cutoff bounds give
$t\le\mathcal B_R(x)u$ and
$\operatorname{widen}_x(t)\le b_x\mathcal B_R(x)u+u$; both exponential
tails are at most $x^{-1}$ eventually, since $B=\lceil(\log x)^2\rceil$ and
$\operatorname{till}(x)$ is a fixed positive fraction of $B$ up to an
additive constant. Thus fixed
polynomials in $1+B$ are subpower in $x$.  The
remaining factors are bounded by the explicit polynomial
\[
 Y(x)=N_0D_xM_0A_0\mathcal B_R(x)
       \bigl(b_x\mathcal B_R(x)+1\bigr)
       +E_0X\mathcal B_R(x).
\]
The source scale bounds give $D_x\le Cb_x^4$, so the displayed
polynomial satisfies $Y(x)\le C'(1+B)^{15}\le x^{1/16}$ eventually.
Hence the right side of
\eqref{eq:ex-lattice-focused-mean-bound} is at most
\[
 Y(x)u^2\le x^{1/16}x^{-1+1/16}=x^{-7/8},
\]
which proves \eqref{eq:ex-lattice-reader-mean}.  These are exactly
the price and widening estimates established in the local Gaussian
source; $B$ remains the outer scale throughout.

\emph{The inverse-column comparison.}
The official-address map preserves physical anchors, row labels and
the sector-membership test.  Its restriction to a panel is onto: a
bare address in the panel is within outer distance $B$ of the cleared
center, so the seed-dome exclusion makes it live in the actual store.
Injectivity is immediate from the official map.  This proves the
finite index bijection only; it does not identify matrix entries.

For each matched pair of indices, the local row estimate gives an
entrywise Gram bound.  For each row constant $A\ge0$, define
\[
 w_x(A,t)=A(1+\operatorname{till}(x))t,\qquad
 \mathfrak m_x(A,t)=D_x\bigl((A+w_x(A,t))w_x(A,t)
                                  +w_x(A,t)A\bigr).
\]
Set $w_x=w_x(A_1,t)$ and $\mathfrak m_x=\mathfrak m_x(A_1,t)$;
then the entrywise error is at most $\delta_x=\mathfrak m_x$.
The moving row differs by at most $w_x$ and the identity row is at
most $A_1$.  Thus, term by term,
\[
 |a_{rj}a_{rk}-b_{rj}b_{rk}|
 \le |a_{rj}|\,|a_{rk}-b_{rk}|
      +|a_{rj}-b_{rj}|\,|b_{rk}|
 \le (A_1+w_x)w_x+w_xA_1.
\]
At most $D_x$ rows contribute, proving the claimed entrywise bound
$\delta_x$.  Let $n_p$ bound the number of coordinates in either
panel.  For positive panel matrices with lower bounds $\lambda$ and
 $\Lambda$, the ordinary resolvent identity gives
\[
 X^{-1}-Y^{-1}=X^{-1}(Y-X)Y^{-1},\qquad
 |(X^{-1}-Y^{-1})_{ij}|
       \le n_p^2\lambda^{-1}\delta_x\Lambda^{-1}.
\]
This is precisely the finite-window comparison; its two distinct
spectral floors are both used.  The count $n_p$ is bounded by the
number of colors times the local jet-slot count at radius $p$.

The nonnegative panel charge is
\[
 \mathcal C_x=n_p^2\lambda^{-1}\mathfrak m_x\Lambda^{-1}.
\]
The local stencil count gives $n_p\le Cb_x^4$ and $D_x\le Cb_x^4$.
Also $t\le Cb_x^5\vartheta(x)$ and
$w_x\le Cb_x^6\vartheta(x)$.  Since $\vartheta(x)\le1$ eventually,
the definition of $\mathfrak m_x$ yields
$\mathfrak m_x\le Cb_x^{16}\vartheta(x)$ and hence
$\mathcal C_x\le Cb_x^{24}\vartheta(x)$.  The subpower scale bound
implies $Cb_x^{24}\le x^{1/8}$ for all sufficiently large $x$; the
temperature bound is $\vartheta(x)\le x^{-1/2+1/8}=x^{-3/8}$.
Their product is $x^{-1/4}$, proving
\eqref{eq:ex-lattice-reader-column}. The source price is $R\chi(x)$;
its fixed-multiplier inequality
$\mathfrak m_x(A,R\chi(x))\le\mathfrak m_x(RA,\chi(x))$
absorbs $R$ into the row constant before applying the source fare
bound. This completes the panel comparison without treating the
curved and bare Gram matrices as equal.
\end{proof}
